\chapter{Orientações sobre o Código de Ética e legislação profissional}
\label{chap:of-codigo_de_etica}

\section{O que é?}

A Comissão de Orientação e Fiscalização (COF) do CRP~SP, descentralizada nas 11 subsedes, oferece à categoria e à sociedade orientação sobre ética e atuação profissional e temas relacionados à legislação do Sistema Conselhos de Psicologia.

\section{Quem pode utilizar esse serviço?}

O acesso é disponibilizado à categoria e à sociedade, não se restringindo apenas às/aos cidadãs/ãos brasileiras/os.

\section{Como solicitar o serviço?}

O serviço é disponibilizado por telefone ou \emph{e-mail} e, presencialmente, pelos\href{http://www.crpsp.org/noticia/view/2460/contatos-das-subsedes-do-crp-sp}{contatosdiretos das subsedes}, e direcionado, de acordo com a demanda e o município, para o território de referência.

O horário de atendimento (presencial e telefônico) da COF é das 10 às 16 horas, de segunda a sexta-feira. O atendimento é feito em todas as 11 subsedes do CRP~SP. Para atendimento presencial, recomendamos entrar em contato com a subsede antecipadamente.

\section{Que informações ou documentos são necessários?}

\subsection{Se o pedido é feito em umasubsede}

\begin{itemize}
  \item  A pessoa que solicita a informação deve se identificar, e os dados são  enviados ao Tribunal de Contas da União para controle (não é gerado  processo administrativo ou legal);
  \item  caso a/o solicitante seja psicóloga/o, também deverá informar o número de inscrição no CRP.
\end{itemize}

\chapter{Se o pedido é feito por \emph{e-mail}}

\begin{itemize}
  \item  Deve-se informar um \emph{e-mail} válido para envio da resposta.
\end{itemize}

\section{Quais são as etapas para a realização desse serviço?}

\subsection{Para orientações solicitadas via \emph{e-mail} da subsede}

\begin{enumerate}
  \item  O pedido de orientação é feito por \emph{e-mail};
  \item  A resposta é enviada (também por \emph{e-mail}).
\end{enumerate}

\subsection{Para orientação presencial}

\begin{enumerate}
  \item  A pessoa que solicita a informação entra em contato com a subsede de  referência para avaliar o melhor horário para o atendimento;
  \item  a informação é prestada na subsede.
\end{enumerate}

\section{Quanto custa?}

Gratuito.

\section{Quanto tempo leva?}

Não há prazo definido. Os atendimentos telefônicos e presenciais são realizados imediatamente pela equipe técnica, salvo quando esta se encontrar em atividade externa. Nesse caso, é oferecida à/ao solicitante a possibilidade de telefonar para outro território e pedir para a/o técnica/o da subsede atender a demanda.

Sendo de interesse falar diretamente com a/o técnica/o de referência, solicita-se o envio de \emph{e-mail} e/ou contato telefônico posterior. Para respostas por \emph{e-mail}, a COF leva em média de 1 a 7 dias para retorno, a depender da complexidade da orientação solicitada.

\section{Legislação relacionada}

\begin{itemize}
  \item  \href{http://www.planalto.gov.br/ccivil_03/leis/l5766.htm}{Lei Federal  nº~5.766/1971}, que cria o Conselho Federal e os Conselhos Regionais  de Psicologia;
  \item  \href{https://atosoficiais.com.br/cfp/resolucao-do-exercicio-profissional-n-10-2005-aprova-o-codigo-de-etica-profissional-do-psicologo}{Resolução  CFP nº~10/2005}, que institui o  \href{https://www.crpsp.org/uploads/pagina/289379/2j9LIMPLJ9jFr5YrK57HmAiBWjVMxdbe.pdf}{Código  de Ética Profissional da/o Psicóloga/o};
  \item  \href{https://transparencia.cfp.org.br/wp-content/uploads/sites/10/2017/07/Resolu\%C3\%A7\%C3\%A3o-CFP-010-2017-Politica-de-Orienta\%C3\%A7\%C3\%A3o-e-Fiscaliza\%C3\%A7\%C3\%A3o-dos-Conselhos.pdf}{Resolução  CFP nº~10/2017}, que institui a Política de Orientação e Fiscalização  do Sistema Conselhos de Psicologia
\end{itemize}

\sumario